\documentclass[11pt]{article}
\usepackage[a4paper,margin=19mm]{geometry}
\usepackage{fontspec}
\setmainfont{lmroman10-regular.otf}
\usepackage{amsmath,amssymb,amsthm,mathtools,amscd,xcolor,graphicx,hyperref}
\hypersetup{hidelinks}
\numberwithin{equation}{section}
\newtheorem{theorem}{Theorem}[section]
\newtheorem{lemma}[theorem]{Lemma}
\newtheorem{proposition}[theorem]{Proposition}
\newtheorem{corollary}[theorem]{Corollary}
\newtheorem{claim}[theorem]{Claim}
\newtheorem{fact}[theorem]{Fact}
\newtheorem{conjecture}[theorem]{Conjecture}
\theoremstyle{definition}
\newtheorem{definition}[theorem]{Definition}
\newtheorem{example}[theorem]{Example}
\newtheorem{remark}[theorem]{Remark}
\newtheorem{question}[theorem]{Question}
\theoremstyle{plain}
\newcommand{\C}{\mathbb C}
\newcommand{\R}{\mathbb R}
\newcommand{\Q}{\mathbb Q}
\newcommand{\Z}{\mathbb Z}
\newcommand{\N}{\mathbb N}
\newcommand{\eee}{\mathrm e}
\newcommand{\cclass}[1]{\textsf{#1}}
\newcommand{\cproblem}[1]{\textsc{#1}}
\newcommand{\ie}{i.\,e.}
\newcommand{\Ie}{I.\,e.}
\newcommand{\eg}{e.\,g.}
\newcommand{\Eg}{E.\,g.}
\newcommand{\phd}{Ph.\,D.}
\newcommand{\msc}{M.\,S.}
\newcommand{\bsc}{B.\,S.}
\newcommand{\expref}[2]{\hyperref[#2]{#1~\ref{#2}}}
\newcommand{\expeqref}[2]{\hyperref[#2]{#1~\eqref{#2}}}
\newcommand{\exprefsilent}[2]{\hyperref[#2]{#1}}
\providecommand{\etal}{et~al.}
\usepackage{eprint}
\newcommand{\TT}{{\bf T}}
\newcommand{\FF}{{\cal F}}
\newcommand{\eps}{{\varepsilon}}
\newcommand{\ignore}[1]{}
\title{Spectral expander reports allow linear-time large-fraction corruption detection}
\author{Noga Alon, Elchanan Mossel, Robin Pemantle}
\date{doi:10.4086/toc.2020.v016a001; CC BY 3.0}
\begin{document}
\maketitle
\noindent\textit{Editorial note.} Complete original mathematical paper, all original author core proofs and published bibliography retained. Original native frontmatter and biographies remain in raw source and publication PDF assets. Existing article interface renders ToC notation for XeLaTeX only.\par\medskip
\noindent\textit{Editorial note.} One general main object only; other source-local families, repeated statements, special cases and direct application corollaries remain complete context. Exact original CC BY3.0 and all native conditions retained.\par\medskip
\noindent\textit{Editorial note.} Authors own mathematics. No proof review, mathematical repair, new assumptions, counterexample search, content hashes or cross-version comparison.\par\medskip
\noindent\textit{Allowed standard prerequisites (editorial).}\begin{itemize}
\item Original graph, circuit or algebraic conventions and cited standard tools. All native model, parameter and field assumptions retained with cited background and complete source-local constructions.. No core author argument removed or replaced by a collector proof; authors retain mathematical responsibility.
\end{itemize}\medskip
\section{Introduction}
We study the problem of 
%
%
 \emph{corruption detection} in networks. %
Given a network of agents, a subset of whom are 
corrupt and the rest are truthful, our goal is to %
find as many truthful and corrupt agents as possible. 
Neighboring agents audit each other. We assume that 
truthful  agents report the %
%
type  %
of their neighbors accurately. 
We make no assumption on the report of corrupt agents. 
For example, two corrupt neighbors may collude and report each other as
non-corrupt. Similarly, a corrupt agent may 
prefer to report the status of some of 
its neighbors accurately, hoping that this will establish a 
truthful record for itself. 
%
%
%
We permit that the corrupt agents    %
coordinate their actions in an arbitrary fashion. 

The corruption model studied here is identical 
to the model of diagnosable systems 
that was introduced by Perparata, Metze and Chien~\cite{PMC67} 
as a model of a digital system 
with many components that can 
%
fail. 
It is assumed that components can test 
some other components. The goal in~\cite{PMC67} and
%
in follow-up  %
work,   %
including~\cite{HA74,KTA75,KR80},  %
%
is to characterize networks that can detect 
%
%
%
%
%
all  %
corrupt  %
nodes 
as long as their number does not exceed a given value.
%
%
%
%
Similar models
%
have been   %
introduced and studied in other areas of computer science, 
including Byzantine computing~\cite{LSP} and intrusion 
detection in the security community~\cite{mukherjee1994network}. 
See also the survey~\cite{lgorzata2004survey}.
A byproduct of this line of research is the VLSI chips puzzle discussed
in~\cite[Problems~4-6]{CLRS}. In this puzzle there are $n$ supposedly
identical VLSI chips that in principle are capable of
testing each other. A basic test involves two chips,
each chip tests the other and reports whether
it is good or bad. A good chip always
%
%
gives an accurate report,   %
but the answer of a bad chip cannot be trusted.
The objective is to find a good chip, or all good chips, assuming
more than half of the chips are good, using the minimum possible number
of tests.  This is (an adaptive version of) 
the corruption detection problem
on a complete graph. 
%
%
%

%
%

The original motivation for our work has been 
corruption detection in social and economic networks,
where the main objective is to understand the
structure of networks that enable one to identify as many of the corrupt
nodes and as many of the truthful ones as possible. 
%
%
%
%
%
%
%
Our goal is to understand theoretically 
which network structures are more amenable to corruption and which are
more robust against it. In particular, is it possible to design sparse 
networks that allow to detect corruption even if there are many corrupt nodes? 
Social scientists have studied many 
%
aspects of corruption in networks, %
see, \eg, %
\cite{nielsen2003corruption,rock2004comparative,fjeldstad2003fighting}.
However, to the best of our knowledge, prior to this
%
%
%
%
work, the effect of the network structure on corruption detection %
has not been systematically studied.     %

As pointed out in
%
follow-up    %
work~\cite{JiMoRa:18}, while 
%
%
%
%
%
%
%
%
``corruption \emph{is} prevalent in many real-world networks, 
%
[\dots]    %
in many scenarios it is not easy to pinpoint even a single 
truthful node'' in spite of our positive theoretical results.
``One reason for that is that some of the 
assumptions do not seem to hold in some 
%
real-world     %
networks. For example, we assume that audits from the truthful nodes are not only non-malicious, but also perfectly reliable. In practice this assumption is unlikely to be true: many truthful nodes could be non-malicious but simply unable to audit their neighbors accurately. Further assumptions that may not hold in some scenarios include the notion of a central agency that is both uncorrupted and has access to reports from every
%
agent      %
and possibly even the assumption that the number of corrupt nodes is less than $|V| / 2$.'' 
In addition there are many constraints on the network that may prohibit it to be an expander. 
Despite these shortcomings of our model, our work points to ideal conditions that allow corruption detection. 
While it is perhaps unreasonable to assume that one imposes an expander auditing structure among government agencies, imposing such a structure 
among more equal entities
%
%
may be more realistic. Furthermore, our results for directed graphs suggest
such structures even in cases where the auditing relation is not symmetric.  


\subsection{Formal definitions and main results} 

\begin{definition}[Digraph, Oriented Graph, Graph] 
For a set $V$, let $V^{(2)}$ denote the set of ordered pairs of 
distinct elements. 
A \emph{digraph} $G=(V,E)$ consists of a set $V$ of nodes (``agents'') 
and a set $E \subset  V^{(2)}$ of directed edges. If $E$ contains no
anti-parallel directed edges (a pair of edges $(u,v)$ and $(v,u)$), then
$G$ is an \emph{oriented graph}. If $E$ contains a directed edge 
$(u,v)$ if and only if it contains $(v,u)$, $G$ is an 
\emph{undirected graph} or simply a \emph{graph}.
%
\end{definition}
%
%
\begin{definition}[Type, Report]
Consider a
%
digraph    %
$G=(V,E)$ along with a function $\tau : V \to \{\mathsf{c},\mathsf{t}\}$ 
that assigns each node a \emph{type}. We call $\tau$ a 
\emph{truthfulness assignment} to $G$. 
We call a node $v$ \emph{truthful} if $\tau(v) = \mathsf{t}$ and corrupt if $\tau(v) = \mathsf{c}$. 
We write $T = \tau^{-1}(\mathsf{t})$ for the set of truthful nodes and $B =  \tau^{-1}(\mathsf{c})$
for the set of corrupt nodes, so that $V = T \sqcup B$ is a partition of the set of nodes. 
A \emph{report} is a function $\rho: E \to \{\mathsf{c},\mathsf{t}\}$. 
A report $\rho: E \to \{\mathsf{c},\mathsf{t}\}$ is \emph{compatible} with 
the truthfulness assignment $\tau$ if for each truthful node $u \in T$ and each
directed edge $(u, v) \in E$ we have $\rho(u,v) = \tau(v)$.  
We call $\rho(u,v)$ the type of $v$ reported by $u$. 
We call $\tau$ a \emph{valid} truthfulness assignment if $|T| > |B|$. 
We say that a report $\rho$ is \emph{feasible} if
it is compatible with at least one valid truthfulness assignment.
%
%
%
\end{definition}
Since it is common to use the letter $C$ for constants, we prefer to denote the set of corrupt nodes by $B$ (the set of  ``bad'' nodes). 

The question we address is, %
under what conditions on the digraph
$G$ and the number of truthful nodes it is possible to determine 
the truthfulness status of
most nodes %
%
with certainty. %
It is easy to see
that this is impossible if $|T| \leq |B|$. Indeed, if $V=V_1 \cup
V_2 \cup W$ is a partition of $V$ into $3$ pairwise disjoint sets
where $|V_1|=|V_2|$ (and $W$ may be empty),
then the corrupt nodes can ensure that all the reports 
in the two scenarios $T=V_1$, $B=V_2 \cup W$ and $T=V_2$, $B=V_1 \cup W$
will be identical. As there is no common truthful node  in these
two possibilities, no algorithm can locate a truthful
node with no error.

Our main result is that if the graph is a good bounded-degree 
directed expander, in the sense described below,
and we have a majority of truthful nodes, then it is possible to
%
determine  %
the truthfulness status of
most of the nodes.   %
%
%

We recall the
%
definition of regular spectral expander graphs.
%
\begin{definition} \label{dec:ndl}
A graph $G$ is an
\emph{$(n,d, \lambda)$-graph} if it is $d$-regular, has $n$ nodes and all
its eigenvalues besides the top one are in absolute value at most
$\lambda$.  Such a graph with $\lambda=2 \sqrt {d-1}$ is called a
\emph{Ramanujan graph}.
\end{definition} 
%
%
%
Explicit constructions of Ramanujan graphs were given by Lubotzky,
Phillips and Sarnak \cite{LPS} and by Margulis \cite{Ma}.
%
%
For the known connection between the 
expansion and pseudo-random properties of graphs and their eigenvalues,
see, \eg, \cite{AC}, \cite{HLW} and the references therein.

%
%
%
%
%
%
%
%
%

%
%
%
%
%
%
%
%
%
%

%
%
%
%
%
%
%
First we consider the somewhat simpler case of \emph{undirected graphs}.
%
%

%

%
The main result for the undirected case
is the following. A more general version is stated as Theorem 
\ref{t21} in subsection \ref{subsec:undirected}.

%
%
%
\begin{theorem}  
\label{t11}
Let $G = (V,E)$ be a Ramanujan graph, that is, an
$(n,d,\lambda)$ graph with $\lambda =2 \sqrt{d-1}$, and
suppose $d \geq 100$.  Then, for every report
$\rho$, there exist sets $T'$ and $B'$ of nodes
such that for every valid truthfulness assignment
$\tau$ compatible with $\rho$ we have
\begin{equation}
\label{eq:main_cond_ram}
T' \subset T, \quad B' \subset B, \quad
|T \setminus T'| \leq \frac{32}{d} |B| , \quad
|B \setminus B'| \leq  \frac{32}{d} |B|,
\end{equation}
where $T = \tau^{-1}(\mathsf{t})$ and $B =  \tau^{-1}(\mathsf{c})$.
%
%
%
%
%
%
Moreover, there is a linear-time
algorithm that finds subsets $T'$ and 
$B'$ satisfying~(\ref{eq:main_cond_ram})
for every truthfulness assignment $\tau$  
compatible with $\rho$ which satisfies
$|T| >  (1 +  {12}/{d}) ({n}/{2})$.
\end{theorem}
Recall that $|T| > |B|$ in the theorem  above since $\tau$ is valid.
%
%
%
%
We make a few remarks on the conclusions of the theorem.
%
%

\begin{enumerate}
\item Given any graph $G$ with all degrees bounded by $d$, with 
at least $2d+1$ corrupt nodes, then there is a set of 
 $\lfloor |B|/(2d+1) \rfloor$ corrupt nodes and a set of 
$\lfloor |B|/(2d+1) \rfloor$ truthful nodes
%
that  %
cannot be %
classified   %
even if in addition to the report $\rho$, we are given the number of 
corrupt nodes. Thus, the
%
fraction of each type of nodes 
we fail to classify is 
%
%
%
%
%
%
tight up to a constant factor. 
To see that this is the case, consider the following selection of $b$ 
%
%
%
corrupt nodes. Choose   %
an arbitrary node $v_1$ in the graph and set all its neighbors $N(v_1)$ to be corrupt, then choose $v_2 \notin \{ v_1 \} \cup N(v_1)$ and set all of its neighbors to be corrupt, then choose one of $v_1,v_2$ to be corrupt 
and the other to be truthful. Continue in a similar fashion, defining $N(v_3),N(v_4)$ to be corrupt and choosing one of $v_3,v_4$ to be corrupt and the other truthful. 
Continue until the number of corrupt nodes left, 
$b'$,   %
satisfies $b' < 2d+1$. 
%
%
Set $b'$ of the remaining nodes to be corrupt. 
%
 Let $r(v,u) = \mathsf{c}$, whenever $v$ is corrupt.  It is then clear that it is impossible to decide which of the $v_i$ is corrupt and which is truthful.
 %
 
\item In the case where $|B|/d = o(n)$,  the theorem allows to recover 
$1-o(1)$ fraction of the good nodes. 

\item 
%
%
%
%
The proof of the theorem is algorithmic.  However,
the algorithm takes exponential time   %
if we only assume that $|T|>|B|$ (or if we assume that
$|T|>(1/2+\mu)n$ for a very small fixed $\mu=\mu(\lambda,d)>0$). 
\end{enumerate} 

The fact that the detection algorithm is not efficient
even in cases when $T$ is larger than $B$ by a (not too large) 
linear number of nodes
%
%
is not a coincidence. Indeed, the algorithm described
in the proof of the theorem, presented
in the next sections, provides a set $T$ of more than $n/2$ truthful
nodes, which is consistent with the reports obtained, when such a
set exists. We show that the problem of producing
such a set when it exists is $\cclass{NP}$-hard,
even when restricted to bounded-degree expanders (and even if 
we ensure that there is such a set of size at least
$n/2+\eta n$).
\begin{theorem}
\label{t12}
%
There exist $c > 0, \eta > 0$ 
such that the following promise problem is 
$\cclass{NP}$-hard. The input is a graph $G=(V,E)$ with $|V| = n$, 
which is an $(n,d,c\sqrt{d})$-graph  
along with a report $\rho$. 
%
%
The promise is that either
\begin{itemize}
\item 
$\rho$ is compatible with $\tau$ satisfying  
$|T| = |\tau^{-1}(\mathsf{t})| \geq n/2 + \eta n$, or
\item 
any $\tau$ which is compatible with $\rho$ satisfies 
 $|T| = \tau^{-1}(\mathsf{t})  \leq n/2 - \eta n$. 
\end{itemize}
The objective is to distinguish between the two options above.
\end{theorem}
The proof is presented in \expref{Section}{subsec:hardness}.
%

We also establish in \expref{Section}{subsec:bad_expanders} the 
following simple statement, 
which shows that at least some
(weak) form of expansion is needed for solving the
corruption detection problem.
\begin{proposition}
\label{p13}
Let $G=(V,E)$ be a graph on $n$ nodes so that it is possible to
remove at most $\epsilon n$ nodes of $G$  and get a graph in which 
%
each   %
connected component is of size at most $\epsilon n$. Then given a report 
$\rho$ compatible with an assignment $\tau$, with $T = \tau^{-1}(\mathsf{t})$ and $B =  \tau^{-1}(\mathsf{c})$,
and $|T| \geq (1-2 \epsilon) n$, it is impossible to
identify even a single member $t \in T$ from the reports of
all nodes. In particular, this is the case for planar graphs or
graphs with a fixed excluded minor even if
$\epsilon=\Theta(n^{-1/3})$.
\end{proposition}

Note that there is still a significant gap between the expansion properties
that suffice for solving the detection problem, described in 
\expref{Theorem}{t11},  and the conditions in the last proposition that are
necessary for such a solution. It will be interesting to obtain tighter
relations between expansion and corruption detection. This is
further discussed in \expref{Section}{sec:open}.


\subsection{Results for directed graphs}
In this subsection we consider directed graphs (digraphs). 
This is motivated by the fact
that in various auditing situations it is 
unnatural to allow $u$ to inspect $v$ whenever 
$v$ inspects $u$. In fact, it may even be desirable not to allow any short cycles in the directed inspection graph. 

%

%

%
For a set of nodes $A$, in a graph $G$, define:
\begin{equation} \label{def:na}
N(A) := \{ v : \exists u \in A, \{u,v\} \in E \}.
\end{equation}

For a digraph $G=(V,E)$, let 
\[
N^{+}(U) = \{ (v : \exists u \in U, (u,v) \in E\}, \quad 
N^{-}(U) = \{ (v : \exists u \in U, (v,u) \in E\}.
\]
Note that $N(U), N^{+}(U)$ and $N^{-}(U)$ may contain elements from $U$. 
We first present an explicit construction of regular directed 
expanders which expand at three 
different scales as in the undirected case. 

%
\begin{proposition}
\label{p14} 
There
%
exist constants $c_1,c_2 > 0$ 
and  infinitely many values of 
$d$ for which there are infinitely many values of $n$ and an 
explicit construction of $3(d-6)$-regular oriented
%
graphs $G=([n],E)$ which satisfy
%
the following properties.
%
\begin{itemize}
\item The girth of $G$ as an undirected graph is at least $2 \log_{d}(n) / 9$. 
\item If $A \subset [n]$ is of size at most $n/(9 d)$ then $|N^{+}(A)| > c_1 d |A|$ and $|N^{-}(A)| > c_1 d |A|$. 
\item For any set $A \subset [n]$ of size at most $n/4$, it holds that $|N^{+}(A)| > 2 |A|$ and $|N^{-}(A)| > 2|A|$. 
\item If $A,B \subset [n]$ with $|A| \geq c_2 n /d$, and
  $|B| \geq n/4$, then there is a directed edge from $A$ to $B$ and a
  directed edge from $B$ to $A$.
\end{itemize} 
\end{proposition} 

%

The proof, presented in \expref{Section}{sec:directed}, is based 
on ``packing'' three different group-based expanders 
to obtain the expansion in the different scales.  Given 
the directed expanders constructed in \expref{Proposition}{p14}, we prove the following directed analogue of \expref{Theorem}{t11}.

\begin{theorem}
\label{t15}
%
Consider the oriented graphs constructed in \expref{Proposition}{p14} . 
There exist
%
%
%
constants $c_3(c_1,c_2)$ and $c_4(c_1,c_2)$,
%
for which the following holds. 

%
For every report
$\rho$, there exist sets $T'$ and $B'$ of nodes
such that for every valid truthfulness assignment
$\tau$ compatible with $\rho$ we have
\begin{equation} \label{eq:main_cond_dir}
T' \subset T, \quad B' \subset B, \quad |T \setminus T'| \leq \frac{c_3}{d} |B| , \quad 
|B \setminus B'| \leq  \frac{c_3}{d} |B|. 
\end{equation}
where $T = \tau^{-1}(\mathsf{t})$ and $B =  \tau^{-1}(\mathsf{c})$.

%

Moreover, there is a linear-time
algorithm that finds subsets $T'$ and
$B'$ satisfying~(\ref{eq:main_cond_dir})
for every truthfulness assignment $\tau$
compatible with $\rho$ which
satisfies
$|T| >  (1 +  {c_4}/{d}) ({n}/{2})$.
%
\end{theorem}

%

\subsection{Relation to distributed computing} 
%
The model considered in the paper requires a central agency that 
obtains the report $\rho$ from all nodes. This is in contrast to models in 
distributed computing and byzantine agreement where all nodes 
only communicate with neighboring nodes. We note, however, that if the 
only objective is to ensure that most %
%
truthful nodes will know the correct
types of large subsets of the set of truthful and the set of
corrupt nodes, this can
be achieved in the common distributed model, even if the number
of truthful nodes is much less than the number of corrupt ones. This
is stated in the next proposition. Its proof, which is much simpler
than that of Theorem \ref{t11}, is presented at
the end of subsection \ref{subsec:undirected}.
%

%
\begin{proposition} 
\label{p:distributed}
Let $G=(V,E)$ be an $(n,d, \lambda)$-graph and suppose that each node
of $G$ has a unique name known to all its neighbors.
Let $\tau$ be a truthful assignment with 
$T = \tau^{-1}(\mathsf{t})$ and $B =  \tau^{-1}(\mathsf{c})$,
put $|T|=tn$, and assume that
$t \geq \frac{3 \lambda}{d}$. Then there is 
an efficient  distributed algorithm on the graph $G$  in which
%
all nodes in a subset $T'$ of the set of truthful nodes $T$ output
subsets $T' \subset T$ and $B' \subset B$, classifying correctly the
type of all nodes in these subsets,  so that the following
holds.  
 %
%
%
\begin{equation} 
\label{eq101}
\mbox{If}~~t >1/2~~ \mbox{then} \quad 
|T \setminus T'| \leq \frac{8\lambda^2}{d^2} |B| , \quad 
|B \setminus B'| \leq  \frac{8 \lambda^2}{d^2} |B|,
\end{equation}
\begin{equation} 
\label{eq102}
\mbox{If}~~\frac{3 \lambda}{d} \leq t \leq 1/2~~ \mbox{then} \quad 
|V \setminus (T' \cup B')| \leq \frac{3\lambda^2}{td^2}n
 ~(~\leq  \frac{\lambda}{d} n).
\end{equation}
%
%
%
%
%
%
%
%
%
%
``Efficient'' here means that every node sends and receives $O(n)$
messages. %
\end{proposition} 

\subsection{Related work} 

The vast literature on corruption detection in computer science,
and in particular on the diagnosable system problem and the PMC model
introduced in~\cite{PMC67},
%
deals  %
either with the problem of identifying all
corrupt nodes, or with that of identifying a single corrupt node. 
As observed in~\cite{PMC67},  a necessary condition 
for the identification of \emph{all} corrupt nodes in a network with $t$ 
corrupt nodes is that the \emph{minimal in-degree} 
in the network is at least $t$.
Therefore, if the number of corrupt nodes is linear in the total number
of nodes, all in-degrees have to be linear, and the total
number of edges has to be quadratic.

The main contribution of the present paper is 
a proof that the number of required edges may be much smaller
when relaxing the requirement of 
identifying \emph{all} corrupt nodes and replacing it by the 
requirement of the 
identification of the types
%
%
of most of the nodes.  %
By relaxing the requirement as above 
we are able to study bounded-degree graphs, in particular $d$-regular graphs. 
Our main new result is that a linear number
%
of edges ensures the recovery of the types of a large fraction
of the nodes, provided the graph is a sufficiently strong expander.
It was shown already
in~\cite{PMC67} that   a linear number of edges suffices
to ensure the detection
of \emph{a single} corrupt node. We show that such a small number of
edges suffices to determine the types of a $1-O(1/d)$ fraction of the 
nodes, even when the number of truthful nodes exceeds that of
corrupt ones by only $1$.

In the context of Byzantine agreement it was discovered already
in the 80s that expanders allow ``almost everywhere agreement.'' This
was first established by Dwork et al.\ in~\cite{DPPU} and  further 
developed in subsequent work, see, in particular~\cite{GO} and its 
references. In~\cite{DPPU} it is shown that on expanders 
one can achieve a relaxed notion of 
Byzantine agreement, where a large fraction of non-faulty nodes agree 
on a value. Like in our results, in the 
%
bounded-degree case,   %
this fraction is bounded away from $1$ 
unless the number of corrupt nodes is sublinear. 
Moreover, the results of~\cite{DPPU} require the number of 
faulty nodes to be a small fraction 
(much smaller than $1/2$) of the total number of nodes. 

It is therefore not very surprising that graph 
expansion is relevant to corruption detection. It is, however, interesting
to see that in sufficiently strong expanders it is possible to classify
most of the truthful and most of the corrupt nodes even if the number of
truthful nodes exceeds the number of corrupt ones by only $1$.

\subsection{Techniques} 

The proof of \expref{Theorem}{t11} rely crucially on the fact that $(n,d,O(\sqrt{d}))$-graphs expand well at different scales.  
%
%
%
%
%
%
\begin{itemize}
\item Every set $A$ of size $\Omega(n/d)$ and every set $B$ of size $n/4$ have at least one edge between them.
\item Every set $A$ of size $O(n/d)$ has $|N(A) \setminus A| \geq |A|$.  
\item Every set such that $|A \cup N(A)| \leq 3n/4$ satisfies $|N(A)| \geq \Omega(d) |A|$. 
\end{itemize} 

%
%
%
%
%
%
%
%

For the directed case, we prove the existence of directed expanders of high girth with analogous properties in 
\expref{Proposition}{p14}.

%
%
%
%
%
%
%
%
%
%
%
%
%
%
%
We observe that a certain weak expansion property, 
namely, the nonexistence of small separators, is 
necessary for corruption detection.
Combining this observation with the 
planar separator theorem of Lipton and Tarjan and its extensions
we conclude that planar graphs and graphs with a fixed excluded minor 
are not good for corruption detection. 

Finally we discuss the algorithmic aspects of our problem 
using results about hardness of approximation.


\section{Proofs} \label{sec:proofs} 

In this section we present the proofs of all results besides the ones
for directed graphs. These appear in the next section.

\subsection{Undirected graphs} 
\label{subsec:undirected}
We first state the following generalization of \expref{Theorem}{t11}. 
\expref{Theorem}{t11} follows when considering graphs 
with $\lambda \leq 2 \sqrt{d-1}  \;\;(< 2 \sqrt d)$. 

\begin{theorem}
\label{t21}

Let $G = (V,E)$ be a $(n,d,\lambda)$ graph, where
$d^2 \geq 24 \lambda^2$.  
Then, for every report
$\rho$, there exist sets $T'$ and $B'$ of nodes
such that for every valid truthfulness assignment
$\tau$ compatible with $\rho$ we have
\begin{equation}
\label{eq:main_cond}
|T \setminus T'| \leq \frac{8 \lambda^2}{d^2} |B| , \quad 
|B \setminus B'| \leq  \frac{8 \lambda^2}{d^2} |B|.
\end{equation}
where $T = \tau^{-1}(\mathsf{t})$ and $B =  \tau^{-1}(\mathsf{c})$.
%
Moreover, there is a linear-time
algorithm that finds subsets $T'$ and
$B'$ satisfying~(\ref{eq:main_cond})
for every
truthfulness assignment $\tau$
compatible with $\rho$ which
satisfies
$|T|>(1/2+3\lambda^2/d^2)n$.
%
\end{theorem}


For a positive $\delta<1/8$
call a graph $G=(V,E)$ on a set of $n$ nodes a
$\delta$-\emph{good}
expander if any set $U$ of at most $2 \delta n$ nodes has 
more than $|U|$ neighbors outside $U$, and there is an edge between
any pair of sets of nodes provided one of them is of
size at least $\delta n$ and the other is of size at least 
$n/4$.
We recall how standard results about expanders 
%
imply that $(n,d,O(\sqrt{d}))$-graphs are $\delta$ good 
for $\delta = \Theta(1/d)$. 
%
%
%
The following lemma is a variant of Tanner's isoperimetric
inequality~\cite{tanner}.
\begin{lemma}[{\cite[Corollary 1]{ABNNR}}]
 \label{cor:expand}
Let $G=(V,E)$ be an $(n,d,\lambda)$-graph and let $B \subset V$ be of size 
$b n$. Let $t n$ denote the size of the set $V \setminus N(B)$
of nodes that have no neighbors in $B$. Then:
\begin{equation}  \label{eq:nonneighbors}
t \leq \frac{\lambda^2}{d^2}\cdot\frac{1-b}{b}.
\end{equation}
\end{lemma}

For completeness, we include the short proof.

\begin{proof}
For $v\in V$ let $N_B(v)=B\cap N(v)$ denote the set of neighbors
of $v$ in $B$.  The Lemma is an immediate consequence of the
following inequality~\cite[Theorem 1]{ABNNR}.
\begin{equation} \label{eq:isoperi}
\sum_{v\in V}\left(|N_B(v)|-bd\right)^2 \le \lambda^2 b(1-b)n\,.
\end{equation}
Indeed, for $v\in V \setminus N(B)$ we have $N_B(v)=\emptyset$
and so summing just for $v\in V \setminus N(B)$ we obtain
from Eq.~\eqref{eq:isoperi} that $tn(bd)^2\le \lambda^2 b(1-b)n$,
which is equivalent to Eq.~\eqref{eq:nonneighbors}.

\noindent
For the proof of Eq.~\eqref{eq:isoperi},
let $A_G$ denote the adjacency matrix of $G$.
Let $f : V\to\R$ denote the vector defined by setting
$$f(v)=\begin{cases} 1-b & \text{ if } v\in B\\
                     -b & \text{ if } v\notin B
        \end{cases} 
$$
So $\sum_{v\in V} f(v)=0$, \ie, $f$ is orthogonal to the 
all-ones vector.  Therefore, $\|A_Gf\|\le \lambda \|f\|$ where
$\|.\|$ stands for the standard euclidean norm.  Easy calculation
shows that $\|A_Gf\|^2$ is equal to the left-hand side of
Eq.~\eqref{eq:isoperi} and $\lambda^2\|f\|^2$ is equal to the 
right-hand side.
\end{proof}
%
%
%
%
%

\begin{lemma} \label{lem:good_expanders} 
Any $(n,d,\lambda)$-graph
in which ${3 \lambda^2}/{d^2} \leq \delta \leq 1/8$, is a 
$\delta$-good expander.
\end{lemma} 

\begin{proof}
Let $U$ be of size at least $\delta n$ and suppose that that there are no edges between 
$U$ and a set $B$ of size $n/4$. Then if $u = |U|/n$, by 
\expref{Lemma}{cor:expand} it follows that 
\[
u \leq \frac{3 \lambda^2}{d^2},
\]
which is a contradiction. 

Similarly, let $U$ be a set of size  $ u n \leq 2 \delta n$ and suppose that 
$|N(U) \setminus U| < u n$. Then the set $B = V \setminus 
(U \cup N(U))$ is of size 
at least $b n$ where $b > 1- 2 u \geq 1/2$. 
This implies by \expref{Lemma}{cor:expand} that
\[
u \leq \frac{\lambda^2}{d^2} \frac{1-b}{b} \leq \frac{\delta}{3} 4 u \leq u/2,
\]
which is a contradiction.
\end{proof} 


The other property of expanders we will need is that small sets expand by a factor of 
$\Omega(d)$. In particular:

\begin{lemma} \label{lem:small_sets} 
Let $G=(V,E)$ be an $(n,d,\lambda)$-graph. 
Let $A \subset V$ be such that $|A \cup N(A)| \leq 3n/4$.
Then 
\begin{equation} \label{eq:small_sets}
| A \cup N(A)| \geq \frac{d^2}{4 \lambda^2} |A|.
\end{equation}
\end{lemma} 

\begin{proof}
The proof is similar. Let $|A| = a n$. 
Let $B = V \setminus (A \cup N(A))$ and put 
$|B| = (1- c a) n \geq n/4$. Since there are no edges between $A$ and $B$, 
it follows that 
\[
a \leq \frac{\lambda^2}{d^2}  \frac{c a}{1 - c a} \leq 4 c a \frac{\lambda^2}{d^2},
\]
and so 
\[
c \geq  \frac{d^2}{4 \lambda^2},
\]
which is the same as~(\ref{eq:small_sets}).
\end{proof}
 



%
%
%
%
%
%
%
%
%
%
%
%
%
%
%


\noindent
%
\begin{proof}[Proof of \expref{Theorem}{t21}]
Let $G$ be an $(n,d,\lambda)$-graph, where
$d^2 \geq 24 \lambda^2$.

Let $\tau$ be a truthfulness assignment consistent with $\rho$ and let $T = \tau^{-1}(\mathsf{t})$ and 
$B = \tau^{-1}(\mathsf{c})$. 
Note that by \expref{Lemma}{lem:good_expanders}, $G$ is $\delta$ good, where 
$\delta = 3\lambda^2/d^2$. 

Let $H$ be the spanning  subgraph of $G$ in which $(u,v) \in E$ iff
$\rho(u,v) = \rho(v,u) = \mathsf{t}$. %
%
Let $V_1,V_2, \ldots ,V_s$ 
be the sets of nodes of the connected  components of 
$H$. 
\begin{claim}
\label{c21}
All the nodes of each $V_i$ are of the same type, that is, for
each $1 \leq i \leq s$, either $V_i \subset T$ or $V_i \subset B$.
\end{claim}

\begin{proof} Suppose $u$ and $v$ are neighbors in $H$. If $u \in T$
then $v \in T$ (as $u$ reports so). If $u \in B$, then 
$v \in B$ (as $v$ reports that $u \in T$).
\end{proof}

Call a component of $H$ truthful if it is a subset of $T$, 
%
otherwise   %
it is a subset of $B$ and we call it corrupt. 

Let $H'$ be the induced subgraph of $G$ on the set $T$ of all 
%
truthful  %
nodes.
\begin{claim}
\label{c22}
Any connected component of $H'$ is also a connected component of 
$H$.
\end{claim}

\begin{proof} If $u,v \in T$ are adjacent in $G$ (and hence in
$H'$), they are
adjacent in $H$ as well, by definition and by the fact that each of
them reports honestly about its neighbors. Thus each component $C'$ of
$H'$ is contained in a component $C$ of $H$. However, no
$v \in T$ is adjacent in $H$ to a node $w \in B$, implying
that in fact $C'=C$ and establishing the
assertion of the claim.
\end{proof}

\begin{claim}
\label{c23}
The graph $H'$ contains a connected component of size at least
$|T|-\delta n > (1/2-\delta)n$.
\end{claim}

\begin{proof} Assume this is false and the largest connected component
  of $H'$ is on a set of nodes $U_1$ of size smaller than
  $|T|-\delta n$. Since the total number of nodes of $H'$ is
  $|T|>n/2$, it is easy to check that one can split the connected
  components of $H'$ into two disjoint sets, each of total size at
  least $\delta n$. However, the bigger among the two is of size
  bigger than $n/4$, and hence, since $G$ is a $\delta$-good expander,
  there is an edge of $G$ between the two groups. This is impossible,
  as it means that there is an edge of $G$ between two distinct
  connected components of $H'$.
\end{proof}

The analysis so far allow us to prove the easy part of the theorem.
\begin{claim}
  If  $|T|>(1/2+\delta)n$ then there exists a linear-time algorithm 
  which finds $T' \subset T$ and $B' \subset B$ such that 
  %
  %
  \[
    |T \setminus T'| \leq 8 \frac{\lambda^2}{d^2} |B| , \quad 
    |B \setminus B'| \leq  8 \frac{\lambda^2}{d^2} |B|. 
  \]
\end{claim}

\begin{proof}
  Note that if $|T|>(1/2+\delta)n$ then \expref{Claim}{c23} implies
  that $H$ must contain a connected component of size bigger than
  $n/2$, which must be truthful. Denote the nodes of this component
  by $T'$ and $B' = N(T')$ and $R = V \setminus (B' \cup T')$. Clearly
  $B' \subset B$ and by \expref{Lemma}{lem:small_sets} it follows that
\[
  |R \cup B'| \geq \frac{d^2}{4 \lambda^2} |R|
\]
which by the assumption that $d^2 \geq 24 \lambda^2$ that gives
$|R|+|B'| \geq 6|R|$ and hence $|R \cup B'| \leq \frac{6}{5}|B'|$
implies
 \[
|T \setminus T'| \leq |R| <
5 \frac{\lambda^2}{d^2} |B| , \quad |B \setminus B'| \leq 
|R|<  5 \frac{\lambda^2}{d^2} |B|, 
\]
establishing the required inequality (with room to spare).


Thus, if this is the case, the set $T'$ is found  
%
by the simple, linear-time algorithm that computes the connected
components of $H$. Furthermore, the set $B'$ can also be found by computing the node boundary of $H$ which is easily computed in linear time as well.
\end{proof}
 
%
%
%
%

It remains to show that even if we only
assume that $|T|>n/2$ then we can still
%
correctly classify   %
most of the truthful and   %
most of the %
corrupt nodes. 
We proceed with the proof of this stronger statement.

By \expref{Claims}{c22} and~\ref{c23} it follows that $H$ contains 
at least one
connected component of size at least $(1/2-\delta)n \geq 3/8 n$.
If $H$ contains only one such component, then this component  
must consist of truthful nodes, and we can
%
correctly classify  %
all of them.
Having these truthful nodes, we also
know the types of all their neighbors. By the assumption on $G$
this gives the types of all nodes 
but less than $\delta n$, thus establishing
(\ref{eq:main_cond}). 
%
%
%
%
%
%
%
%
%
%
%
%
%

 
Otherwise, there is another connected component of size at least
$(1/2-\delta)n$, and as there is no room for more than two such
components, there are exactly two of them, say $V_1$ and $V_2$. Note
that by the expansion properties of $G$ there are edges of $G$ between $V_1$
and $V_2$, and hence it is impossible that both of them are
truthful components.  As one of them must be truthful, it follows
that exactly one of $V_1$ and $V_2$ is a truthful component and the 
other corrupt. We next show that we can determine the types of both
components.  %

Construct an auxiliary weighted graph $S$ on the set of nodes 
$1,2,
\ldots ,s$ representing the connected components  $V_1, V_2
\ldots ,V_s$ as follows. The weight $w_i$ of $i$ is defined
by $w_i= {|V_i|}/{|V|}$. Two nodes $i$ and $j$ are
connected iff there is at least one edge  of $G$ that connects a
node in $V_i$  with one in $V_j$. Call an independent set in
the graph $S$ \emph{large} if its total weight is bigger than $1/2$. Note
that by the discussion above $T$ must be a union of the form
%
$T=\bigcup_{i \in I} V_i$,  %
where $I$ is a large independent set in the
graph $S$. In order to complete the argument we prove the following.
\begin{claim}
\label{c24}
Either there is no 
large independent set  in $S$ containing $1$, or there is no
large independent set in $S$ containing $2$.
\end{claim}

\begin{proof}
  Assume this is false, and let $I_1$ be a large independent set in
  $S$ containing $1$, and $I_2$ a large independent set in $S$
  containing $2$. To get a contradiction we show that for
  $ w(I_1)=\sum_{i \in I_1}w_i$ and $w(I_2) =\sum_{i \in I_2} w_i$ we
  have $w(I_1) + w(I_2) \leq 1$ (and hence it is impossible that each
  of them has total weight bigger than a half).

  To prove the above note, first, that the two nodes $1$ and $2$ of
  $S$ are connected (as each corresponds to a set of more than
  $(1/2-\delta)n$ nodes of $G$, hence there are edges of $G$
  connecting $V_1$ and $V_2$). Therefore $I_1$ must contain $1$ but
  not $2$, and $I_2$ contains $2$ but not $1$.

  If there are any nodes $i$ of $S$ connected in $S$ both to $1$
  and to $2$, then these nodes belong to neither $I_1$ nor $I_2$,
  as these are independent sets. Similarly, if a node $i$ is
  connected to $1$ but not to $2$, then it can belong to $I_2$ but not
  to $I_1$, and the symmetric statement holds for nodes connected
  to $2$ but not to $1$. So far we have discussed only nodes that
  can belong to at most one of the two independent sets $I_1$ and
  $I_2$. If this is the case for all the nodes of $S$, then each of
  them contributes its weight only to one of the two sets and their
  total weight would thus be at most $1$, implying that it cannot be
  that the weight of each of them is bigger than $1/2$, and completing
  the proof of the claim. It thus remains to deal with the nodes of
  $S$ that belong to both $I_1$ and $I_2$.  Let
  $J \subset \{3,4, \ldots ,s\}$ be the set of all these nodes.
  Note, first, that the total weight of the nodes in $J$ is at most
  $2 \delta$, as the total weight of $1$ and $2$ is at least
  $2 (1/2-\delta)=1-2 \delta$.  Note also that by the discussion above
  each $j \in J$ is not a neighbor of $1$ or of $2$.  By the
  assumption about the expander $G$ the total weight of the nodes
  that are neighbors of nodes in $J$ and do not belong to $J$ is
  bigger than the total weight of the nodes in $J$.  Indeed, this
  is the case as the number of neighbors in $G$ of the set
  $\bigcup_{j \in J} V_j$ that do not lie in this set is bigger than the
  size of the set. We thus conclude that if $J'=N_S(J)-J$ denotes the
  set of neighbors of $J$ that do not belong to $J$, then the total
  weight of the nodes in $J'$ exceeds the total weight of the
  nodes in $J$, and the nodes in $J'$ belong to neither $I_1$
  nor $I_2$.  We have thus proved that the sum of weights of the two
  independent sets $I_1$ and $I_2$ satisfies
  \[
    w(I_1)+w(I_2) \leq 2 w(J) +(1-w(J)-w(J')) \leq w(J)+w(J')
    +(1-w(J)-w(J')) =1
  \]
  contradicting the fact that both $I_1$ and $I_2$ are large. This
  completes the proof of the claim.
\end{proof}

By \expref{Claim}{c24} we conclude that one can determine the types of
the components $V_1$ and $V_2$. This means that we can classify
at least $(1/2-\delta)n$ truthful nodes with no error. Recall
that this is
the case also when $H$ has only one connected component of size
at least $(1/2-\delta)n$. Having these truthful nodes, we also
know the types of all their neighbors. By the assumption on $G$
this gives the types of all nodes but less than $\delta n$,
completing the proof of (\ref{eq:main_cond}).   
\vspace{0.1cm}

\noindent
As noted above,  %
the algorithm described in the proof above is a
linear-time algorithm provided $|T|>(1/2+\delta)n$.
However, if we only assume that $|T|>|B|$ the proof
provides only a non-efficient algorithm  for deciding the types of
the components $V_1$ and $V_2$. Indeed, we have to compute the
maximum weight of an independent set containing $1$
in the weighted graph $S$,
and the maximum weight of an independent set 
containing $2$. By the proof above, exactly one of this maxima 
is larger than $1/2$, providing the required types. 
%
\end{proof}    %

\bigskip
The proof of \expref{Proposition}{p:distributed} follows from 
the fact that it deals with the case of a large component 
of truthful nodes that can communicate with each other. 

%
%
%
\begin{proof}[Proof of \expref{Proposition}{p:distributed}]
%
For simplicity, consider a synchronized protocol. In the protocol, 
in each round, each node transmits to all of its
neighbors the  %
%
types  %
of all nodes it recognizes as truthful and as corrupt.  
A truthful node adds to its record 
the types of all nodes it receives from truthful neighbors. Let
$T'$ denote the set of nodes of the 
largest connected component of the induced 
subgraph of $G$ on the set $T$ of truthful nodes and
put $B'=N(T')-T'$. 
It is clear that all nodes in $T'$ following the protocol
classify all elements of $T'$ as truthful and all elements 
of $B'$ as corrupt. 
If $t >1/2$, that is,
$|T| > |B|$, then by the proof of Theorem \ref{t21}, 
(\ref{eq101}) holds.  If $\frac{3\lambda}{d} \leq t \leq 1/2$
then by Lemma \ref{cor:expand} $T'$ is of size  at least, say,
$\frac{tn}{3}$ (since there is at least one edge
between any two sets, each of size at least $\frac{tn}{3} \geq 
\frac{\lambda}{d}n$.)
Applying Lemma \ref{cor:expand} again it follows that
$|V-N(T')| \leq \frac{\lambda^2}{d^2}\frac{3}{t}n \leq
\frac{\lambda}{d} n$ and (\ref{eq102}) follows as all nodes  in
$T'$ classify correctly the types of all nodes in $N(T')$. 
\end{proof} 


\subsection{Hardness} \label{subsec:hardness}

In this subsection we prove \expref{Theorem}{t12} which explains 
the non-efficiency of the algorithm in the proof of \expref{Theorem}{t11}.
\vspace{0.2cm}

\noindent
%
\begin{proof}[Proof of \expref{Theorem}{t12}]
The proof is based on the following result~\cite{BK99}: there exist
constants $b<a<1/2$ such that deciding if a graph $H$ on $m$ nodes,
all of whose degrees are bounded by $4$, has a maximum independent set
of size at least $(a+b) m$ or at most $(a-b) m$ is $\cclass{NP}$-hard.
It is easy to see that in fact one can assume that 
$H$ is $4$-regular. 
 
Let $G'$ be an $(n,d-4,c_1\sqrt{d})$-graph with node-set $V$, where
$c \sqrt{d} \geq 8$.  
Split the nodes into $3$ disjoint sets, $V_1,V_2,V_3$, each
of size $\Omega(n)$, 
where $V_3$ is an independent set in $G'$ of size $m$, all its neighbors
are in $V_2$, $|V_1|=n/2-am$ and $|V_2|=n/2-m+am$, 
where $m = \Omega(n)$. Write $\eta$ for the constant satisfying $b m=\eta n$. 
%
%
Add on $V_3$ a bounded-degree graph $H$ as above, in which it is hard to
decide if the maximum independent set is of size at least $(a+b)m$ or at
most $(a-b)m$. That is, identify the set of nodes of $H$ with $V_3$
and add edges between the nodes of $V_3$ as in $H$. 
Add a $4$-regular graph on $V_1$ and another one on $V_2$. 
Call the resulting
graph $G$ and note that it is an $(n,d,c_1\sqrt{d} + 4)=
(n,d,c\sqrt d)$ graph. 

The reports of the nodes are as follows.  Each node in $V_1$
reports true on each neighbor it has in $V_1$, and corrupt on  any other
neighbor. Similarly, each node of $V_2$ reports true on any neighbor
it has in $V_2$ and corrupt on any other neighbor, and each node in $V_3$
reports corrupt on all its neighbors.  Note that with these reports the
connected components of the graph $H$ in the proof of \expref{Theorem}{t11}
are $V_1$, $V_2$ and every singleton in $V_3$.

It is easy to check that here if $H$ has an independent set $I$ of size
at least $(a+b)m$, then $G$ has a set $T$ of truthful  nodes of size
at least $n/2+bm$, namely, the set $I \cup V_1$, which is consistent with
all reports. If $H$ has no independent set of size bigger than $(a-b)m$,
then $G$ does not admit any set $T$ of truthful nodes of size bigger
than $n/2 - bm$ consistent with all reports. This completes the proof.
%
%
\end{proof}

\subsection{Graphs with small separators} \label{subsec:bad_expanders}

In this subsection we describe the simple proof of \expref{Proposition}{p13}. 
\vspace{0.3cm}

\begin{proof}[Proof of \expref{Proposition}{p13}]
Let $B'$ be a set of at most $\epsilon n$ nodes of $G$ whose
removal splits $G$ into connected components with node-classes
$V_1,V_2, \ldots ,V_s$, each of size at most $\epsilon n$. Consider
the following $s$ possible scenarios $R_i$, for $1 \leq i \leq s$.

\smallskip
\noindent
$R_i$:  the set of corrupt nodes is $B=B' \cup V_i$, all the others
are 
%
truthful   %
nodes. The nodes in $B'$ report that all their
neighbors are corrupt. The nodes in $V_i$ report that their
neighbors in $V_i$ are in $T$, and that all their other neighbors 
are in $B$. (The truthful nodes, of course, report truthfully
about all their neighbors).

\smallskip
\noindent
It is not difficult to check that in all these $s$ scenarios, all
nodes make exactly the same reports. On the other hand, there is
no node of $G$ that is truthful in all these scenarios, hence 
it is impossible to identify a truthful node with no
error. Since the number of corrupt nodes in all scenarios is at most
$2 \epsilon n$, the first assertion of the theorem follows. The
claim regarding planar graphs and graphs with excluded minors
follows from the results in~\cite{LT}, \cite{AST}.
\end{proof}

\section{Directed graphs} \label{sec:directed}

\subsection{Construction of directed expanders}


Here we provide the proofs for the case of directed graphs. We start with
the proof of existence of explicit directed expanders that expand 
in three scales. 
We will use the following lemma for undirected graphs from~\cite{AKRS}. 
\begin{lemma}[Lemma 3.6~\cite{AKRS}]  \label{lem:frac_expand}
Let $G=(V,E)$ be an $(n,d,\lambda)$ graph. Let $6/d < \gamma < 1$. Let $A,X \subset V$ be sets of nodes such that 
\begin{itemize}
\item $|A| \leq \frac{\gamma}{2(d+1)} n$,
\item for all $v \in A$, it holds that $|\{ w \in X : (w,v) \in E\}| \geq \gamma(d+1)$.
\end{itemize}
Then: 
\[
|X| \geq \frac{\gamma^2 d^2}{9 \lambda^2} |A|.
\]
\end{lemma} 

We will also use the following variant of \expref{Lemma}{lem:good_expanders} whose proof is similar. 

\begin{lemma} \label{lem:good_expanders2} 
Any $(n,d,\lambda)$-graph
in which $16 {\lambda^2}/{d^2} \leq \delta$ satisfies that 
%
%
any set of size at least $\delta n / 2$ and any set of size $n/8$ have at least one edge between them.
%
%
%
\end{lemma} 

\begin{proof}
Let $U$ be of size at least $\delta n/2$ and suppose that that there are no edges between 
$U$ and a set $B$ of size $n/8$. Then if $u = |U|/n$, by \expref{Corollary}{cor:expand} it follows that 
\[
u \leq 7 \frac{\lambda^2}{d^2},
\]
which is a contradiction. 

%
%
%
%
%
%
%
%
\end{proof} 




\begin{proof}[Proof of \expref{Proposition}{p14}]
Let $G'=([n],E')$ be a $d$-regular 
undirected non-bipartite Ramanujan Cayley graph as constructed in~\cite{LPS} or~\cite{Ma}. 
%
This is a Cayley graph of a finite group $\Gamma$ of
%
order  <--
$n$, with respect to a set $S'$ of $d$ generators,
$S' = \{a_1, a_1^{-1}, a_2, a_2^{-1},\ldots, a_{d/2}, a_{d/2}^{-1}\}$. 
This graph is known to have girth at least $2 \log_d(n)/3 $, which is
the same as saying there is no nontrivial word of the generators of
length less than $2 \log_d(n)/3$ that is the identity. 

Let
\[
  T = \bigl\{a_4,a_4^{-1},\ldots,a_{d/2},a_{d/2}^{-1}\bigr\}.
\]
Note that the Cayley graph of $\Gamma$ with respect to $T$ is a
%
$(d-6)$-regular  %
graph whose second eigenvalue is at most $2 \sqrt{d} +6$.
Thus if $ d \geq 36$ this is an $(n,d-6,3 \sqrt{d})$-graph.

Now for $1 \leq i \leq 3$, let $T_i = a_i^{-1} T a_i$. 
Then clearly, $G_i$, the Cayley graph  of $\Gamma$ with respect to $T_i$,
is also an $(n,d-6,3 \sqrt{d})$-graph. Moreover,
if $S = T_1 \cup T_2 \cup T_3$, then the Cayley graph $H$ of 
$\Gamma$ with respect to $S$ has girth at least
\[
  \frac{2}{9} \frac{\log n}{\log d},
\]
since a nontrivial word in $S$ of length $k$ corresponds to a
non-trivial word of length at most $3 k$ in $S'$.  Note also that $G'$
is $3(d-6)$-regular.

The desired graph, $G=([n],E)$, will be obtained by assigning directions
to the edges $\bar{E}$ of $H$ as follows: 
%
%
\begin{itemize}
\item If $\{a,b\}$ is an edge of $G_1$ then orient it from $a$ to $b$ if $a > b$ and from $b$ to $a$ if $b > a$.
\item If $\{a,b\}$ is an edge of $G_2$ then orient it from $b$ to $a$ if $a > b$ and from $a$ to $b$ if $b > a$.
\item In the graph $G_3$ all the degrees are even. Pick an orientation
of the edges of 
$G_3$ by picking a directed Eulerian cycle and orienting the edges
according to it. In particular, note that every
in-degree and out-degree is exactly $d/2-3$ in $G_3$. 
\end{itemize} 

We now verify the expansion at the three different scales.
First, we apply \expref{Lemma}{lem:frac_expand} to the graph $G_3$ and
sets $A$ of size $|A| \leq n/(9d)$, 
%
%
%
 and $\gamma = 1/3$ and obtain that 
\[
|N^{+}(A)| \geq \frac{(d-6)^2}{81 (3 {\sqrt d})^2} |A| \geq \frac{d}{400} |A|,
\]
for $d \geq 36$. The proof for $N^{-}(A)$ is identical. 
 

Next, let $A$ and $B$ be two sets of nodes with $n/16 \geq |A| \geq 200 n / d$
and $|B| \geq n/4$. Let $m(A)$ and $m(B)$ denote the medians of the sets of
numbers given by $A$ and $B$.  Then either $m(A) \geq m(B)$ or 
$m(B) \geq m(A)$. If  $m(A) \geq m(B)$, let $A' = \{ v \in A : v \geq m(A)\}$
and $B' = \{v \in B : v \leq m(B)\}$.
Then $|A'| \geq 100 n / d$ and $|B'| \geq n/8$ and all elements of $A'$ are
bigger than all elements of $B'$. 
Thus by \expref{Lemma}{lem:good_expanders2} applied to $G_1$, with
$\delta = 200/d$, 
$\lambda = 3 \sqrt{d}$ and $d-6 \geq 30$, 
there is at least one edge in 
 $G_1$ connecting $A'$ and $B'$. 
This is a directed edge from $A$ to $B$. A symmetric argument applies
if $m(A) \leq m(B)$. 

Note that the last statement implies that if $n/4 \geq |A| \geq 200 n / d$
then the set of non-neighbors of $A$ is of size at most $n/4$ and
therefore $|N(A)| \geq 3 n/4 > 2 |A|$. 

For sets $A$ with $n / (10d) \leq |A| \leq 200 n / d$, we have 
\[
|N(A)| \geq  \frac{d}{400} \frac{n}{10 d} = \frac{n}{4000},
\]
which is greater than $400 n / d$ for $d$ sufficiently large.  
%
\end{proof}
%
%

\noindent
We next present the proof of \expref{Theorem}{t15}, which resembles that of
\expref{Theorem}{t11} but requires several additional ideas.
\vspace{0.2cm}

%
\begin{proof}[Proof of \expref{Theorem}{t15}]   %
Let $c_1,c_2$ be the
constants from \expref{Proposition}{p14} and put
\[
\delta = c_2/d.
\]

Let $\tau$ be a truthfulness assignment consistent with $\rho$ and
let $T = \tau^{-1}(\mathsf{t})$ and 
$B = \tau^{-1}(\mathsf{c})$. 
Let $H$ be the spanning  subgraph of $G$ in which an edge $(u,v)$
of $G$ is an edge of $H$ iff $\rho(u,v) = \mathsf{t}$.
%
%
Let $V_1,V_2, \ldots ,V_s$ 
be the sets of nodes of the strongly connected  components 
(SCCs, for short) of $H$. 
\begin{claim}
\label{c25}
All the nodes of each $V_i$ are of the same type, that is, for
each $1 \leq i \leq s$, either $V_i \subset T$ or $V_i \subset B$.
\end{claim}

\begin{proof} If $u \in T$ and $v$ is an out-neighbor of
$u$ in $H$, then 
$v \in T$ (as $u$ reports so). If $v \in B$, and $u$ is an
in-neighbor of $v$ in $H$, then 
$u \in B$ (as $u$ reports that $u \in T$).
\end{proof}

Call  an SCC  
of $H$ truthful if it is a subset of $T$, else it
is a subset of $B$ and we call it corrupt. 

Let $H'$ be the induced subgraph of $G$ on the set $T$ of all truthful 
nodes.
\begin{claim}
\label{c26}
Any SCC of $H'$ is also an SCC
of $H$.
\end{claim}

\begin{proof} If $u,v \in T$  and $(u,v)$ is an edge of $G$, then
it is an edge of $H$ too.
Thus each SCC $C'$ of
$H'$ is contained in an SCC $C$ of $H$.  This SCC is
truthful, by \expref{Claim}{c25}, and cannot contain any additional 
truthful nodes as otherwise these belong to $C'$ as well. 
\end{proof}

\begin{claim}
\label{c27}
The graph $H'$ contains an SCC of size at least
$|T|-2 \delta n > (1/2-2 \delta)n$. 
%
%
\end{claim}

\begin{proof} Consider the component graph of $H'$: this is the 
directed graph $F$ whose nodes are all the SCCs of 
$H'$, where there is a directed 
edge  from $C$ to $C'$ iff there is some edge of  $H'$ from
some node of $C$ to some node of $C'$.  It is easy and well
known that this graph is a directed acyclic graph, and hence there
is a topological order of it, that is, a numbering
$C_1, C_2, \ldots ,C_r$ of the components  so that all edges
between different components  are of the form $(C_i,C_j)$ with
$i <j$. Order the nodes of $H'$ in a linear order according to
this topological order, where the nodes of $C_1$ come first
(in an arbitrary order),
those of $C_2$ afterwards, etc. Let $u_i$ be the node in place
$i$ according to this order $(1 \leq i \leq |T|)$. If the
nodes $u_{\delta n}$ and $u_{|T|-\delta n+1}$ 
belong to the same SCC, then this component is of size
at least $|T|-2 \delta n$ and we are done. Otherwise, the SCC
containing $u_{|T|/2}$ differs from either that containing
$u_{\delta n}$ or from that containing $u_{|T|-\delta n+1}$.
In the first case, the set $A$ of all SCCs up to that  containing
$u_{\delta n}$ is of size  at least $\delta n$, and the set
$B$ of all SCCs starting from that containing $u_{|T|/2}$ 
is of size at least $|T|/2 \geq n/4$. In addition there is no edge directed
from $B$ to $A$, contradicting the property of $G$. The second case
leads to a symmetric contradiction, establishing the claim.
\end{proof}

Note that the above shows that if $|T|>(1/2+2\delta)n$ then
$H'$ and hence also $H$ must contain a SCC $C$ of size bigger than
$n/2$, which must be truthful.  Let $T'$ denote the set of nodes
reachable from $C$ in $H$ and note that $T' \subset T$ and moreover
$B' := N^{+}(T') \subset B$.
Let  $R = V \setminus (B' \cup T')$. 

By the expansion properties of $G$ it follows that 
$|R| \leq c_2 n / d$ and 
and that $|R| \leq c_3 |B'|/d \leq c_3 |B|/d$ 
for a constant $c_3=c_3(c_1,c_2)$. 

We next show that even if we only
assume that $|T|>n/2$ we can still classify correctly most of the
truthful nodes. 

By the last two claims it follows that $H$ contains at least one
SCC of size at least $(1/2-2 \delta)n \geq 3/8 n$.
If $H$ contains only one such component, then this component  
must consist of truthful nodes, and we can identify all of them
(and hence also the types of all their out-neighbors). 
Otherwise, there is another SCC of size at least
$(1/2-\delta)n$, and as there is no room for more than two such
components, there are exactly two of them, say $V_1$ and $V_2$. Note
that by the properties of $G$ there are edges of $G$ from $V_1$
to $V_2$ and from $V_2$ to $V_1$, 
and hence it is impossible that both of them are
truthful components.  As one of them must be truthful, it follows
that exactly one  of them is truthful  and one is corrupt.
We next show that we can determine the types of both
components.

Recall that we have the SCCs of $H$, and the set $T$ of all truthful
nodes must be a union of a subset of these SCCs.
In addition, this set must be of size bigger than $n/2$ and must be
consistent with all reports of the nodes along every edge (in
the sense that for any edge $(u,v)$ with $u \in T$, the report of
$u$ on $v$ should be consistent with the actual type of $v$).

\begin{claim}
\label{c28}
Given the strongly connected components $V_1,V_2, \ldots ,V_s$
of $H$ and the reports along each edge, either there is no 
union $I_1$ of SCCs including 
$V_1$ whose size exceeds $n/2$ so that $T=I_1,
B=V-I_1$ is
consistent with all reports along the edges, or there is no union
$I_2$ of SCCs including $V_2$ whose size exceeds $n/2$ so that 
$T=I_2,B=V-I_2$ is consistent with all reports along the edges. 
\end{claim}
\vspace{0.1cm}

\begin{proof}
Assume this is false, and let $I_1, I_2$ be as above.
By the above discussion we know that $I_1$ contains
$V_1$ but not $V_2$ and $I_2$ contains $V_2$ but not $V_1$.
Note that if some SCC $V_i$ is contained both in $I_1$ and in $I_2$ 
and there is any directed edge $(u,v)$ from $V_i$ to some other SCC
$V_j$, then if the report along this edge is that $v$ is truthful,
then $V_j$ must be  truthful  component in both $I_1$ and in $I_2$.
Similarly, if the report along this edge is $v \in B$, then $V_j$
must be outside $I_1$ and outside $I_2$. In particular, there are
no edges at all from $V_i$  to  $V_1$ or $V_2$ (as each of them
lies in exactly one of the two unions $I_1$, $I_2$). Let $J$ be the
set of all SCCs that are contained in both $I_1,I_2$. By the remark
above, for every edge $(u,v)$ from a node of $J$ to a node outside $J$,
the report along the edge must be  $v \in B$ (since otherwise $v$
would also be in an SCC which is truthful both in $I_1$ and in
$I_2$ and hence would be in $J$). Thus all edges $(u,v)$ as above
report $v \in B$, implying that all components outside $J$ 
to which there are
directed edges from nodes in $J$  belong to neither $I_1$ nor
$I_2$. By the properties of our graph the total size of these
components exceeds that of $J$, (as $|J| \leq 4 \delta n$ and all
out-neighbors of $J$ are outside $V_1,V_2$), and this shows that
the sum of the sizes of $I_1$ and $I_2$ is at  most
$$
2|J| +(|V|-|J|-|N^+(J)-J|) \leq |V|.
$$
Therefore  it cannot be that both $I_1$ and $I_2$ are of size bigger
than $|V|/2=n/2$, proving the claim.
\end{proof}

By the last claim it follows that one can determine the types of
the SCCs $V_1$ and $V_2$. This means that we can classify
at least $(1/2-2 \delta)n$ truthful nodes with no error. Recall
that this is
the case also when $H$ has only one SCC of size
at least $(1/2-\delta)n$. Having these truthful nodes, we also
know the types of all their out-neighbors. By the assumption on $G$
this gives the types of all nodes but at most $O(|B|/d)$,
completing the proof of the main part of the theorem.  
\vspace{0.1cm}

\noindent
The comment about the
%
linear-time   %
algorithm provided $|T|>(1/2+2\delta)n$
is clear. If we only assume that $|T|>|B|$ the proof
provides only a non-efficient algorithm  for deciding the types of
the SCCs $V_1$ and $V_2$. Indeed, we have to check
all $2^s$ possibilities of the types of each of the SCCs and see
which ones are consistent with all reports and are of total size
bigger than $n/2$.
By the proof above, only one of the two SCCs $V_1,V_2$ will appear 
among the truthful SCCs of such a possibility.
%
\end{proof}   %

\section{Discussion and open problems}
\label{sec:open}

The usefulness of expanders for corruption detection raises 
the natural question about the existence of good explicit 
spectral expanders with any desired number of nodes. 
After our work has been posted, explicit construction for 
such undirected expander graphs with any (large) 
%
number of nodes were 
obtained in~\cite{A19}. 

It is interesting to study in more detail the relation
between expansion and corruption detection.

\begin{question}
Provide sharp criteria in terms of expansion and the fractional 
size of the set $T$ for enabling corruption detection.
\end{question} 

%
%
%
For a weak result in the desired direction,  %
consider the following argument.
We say that an undirected  graph $G$ is \emph{$\delta$-connected}
 if for every two disjoint sets
$A_1, A_2$ with $|A_1| \geq \delta n, | A_2 | \geq (1- 3 \delta) n$
there is at least one edge between $A_1$ and $A_2$. Note that the notion
of
%
$\delta$-connectedness   %
is much weaker than expansion.  In particular a graph
$G$ can be
%
$\delta$-connected,    %
yet at the same time have $\delta n/2 $
isolated nodes, while any $\delta$-good expander must be connected.

\begin{proposition} \label{c:41} 
%
Suppose that $|T| = (1-\epsilon)n$ and the 
graph $G$ is $\epsilon$-connected.
Then, given the reports of all nodes, 
it is possible to find  in linear time a subset
%
%
$T' \subset T$ of size at least $(1-2\epsilon) n$. 
\end{proposition}

\begin{proof}
Let $E' \subset E$ be the set of edges both of 
whose end-points declare each other truthful. 
Recall that each connected component of $G'=(V,E')$ is either truthful or corrupt. 
%
%
%
%

Let $T_1, T_2,\ldots$ denote all the components of size at least $\epsilon
n$ in $G'$. Then we claim that for $T' = \bigcup T_i$, $|T \setminus T'|
< \epsilon n$. Assume otherwise. Since all the connected components of
$T \setminus T'$ are of size at most $\epsilon n$, there exists $T''
\subset T \setminus T'$ of size in $[\epsilon n, 2 \epsilon n]$ with no
edges to 
%
$T \setminus T''$ whose size is in $[(1-3\epsilon)n ,(1-2\epsilon)
n]$. This is a contradiction to $\epsilon$-connectedness,
%
proving the claim. By $\epsilon$-connectedness we cannot have more
than one component $T_i$ of truthful nodes (as there are
edges between any such $T_i$ and the union of the other large
components of truthful nodes), hence $T'$ must be connected. Since it
is clearly easy to find $T'$ 
%
in linear time  given the reports of all nodes,
the desired result follows.
\end{proof}

To see that the conditions of \expref{Claim}{c:41} are tight up to constant
factors, consider the star graph with $m$ leaves.  Assume that 
%
$|T| \leq m-1$. 
Then it is easy to see that one cannot find even one member of $T$
if all nodes declare all their neighbors corrupt. On the other hand,
this example is (vacuously) $1/(4m)$ connected.  To get a non-trivial
example, one can replace each node with a complete graph $K_k$ and each
edge with a complete bipartite graph $K_{k,k}$ for an arbitrary $k$.

In 
%
follow-up    %
work~\cite{JiMoRa:18}, the connection between expansion and corruption detection is formalized using the conjectured hardness of Small Set Expansion~\cite{RS}. Assuming the hardness of Small Set Expansion, 
it is shown that it is computationally hard to approximate the minimal number of nodes whose corruption makes identifying even one truthful node impossible. 

We conclude with a short discussion of a variant of the model. 
{}From the modeling perspective, it is interesting to consider 
probabilistic variants of the corruption detection problem. 
\begin{question}
What is the effect of relaxing the assumption that truthful nodes
always report the status correctly? Suppose for example that each
truthful node reports the status of each of its neighbors independently
accurately with probability $1-\epsilon$.  Note that in this case it is
impossible to detect the status of an individual node with probability
one. However it is still desirable to find sets $T'$ and $B'$ such that
the symmetric difference $T \bigtriangleup T'$ and $B \bigtriangleup B'$
are small with high probability. Under
what conditions can this be achieved? What are good algorithms for
finding $T'$ and $B'$?
\end{question} 

%
%
%
%
\begin{remark}
Alweiss~\cite{Alweiss:19} addressed this problem
after our work has been posted.
\end{remark}     %

\vspace{0.2cm}

%
%
\paragraph{Acknowledgments.}    %
We thank Gireeja Ranade for suggesting to consider
problems of corruption
%
in   %
networks, Peter Winkler for helpful
comments regarding the Byzantine Agreement problem, and two 
anonymous referees for providing relevant references. We are also 
grateful to Laci Babai for 
numerous suggestions to improve the presentation as well as some of
the results. 
%

%
%
%
%
%

%



%% Theory of Computing
%% File: v016a001.bbl
%% Publication date: March 30, 2020
%% Authors: Noga Alon, Elchanan Mossel, and Robin Pemantle
% Define empty bibhead if not already defined
\providecommand{\bibhead}[1]{}
% Define tocrefpdfbookmark if not already defined
\expandafter\ifx\csname pdfbookmark\endcsname\relax%
  \providecommand{\tocrefpdfbookmark}{}
\else\providecommand{\tocrefpdfbookmark}{%
   \hypertarget{tocreferences}{}%
   \pdfbookmark[1]{References}{tocreferences}}%
\fi

\tocrefpdfbookmark
\begin{thebibliography}{10}

\bibitem{A19}\bibhead{A19}
{\sc Noga Alon}: Explicit expanders of every degree and size, 2020.
\newblock [\epfmt{arxiv}{2003.11673}]

\bibitem{ABNNR}\bibhead{ABNNR}
{\sc Noga Alon, Jehoshua Bruck, Joseph~Seffi Naor, Moni Naor, and Ron~M. Roth}:
  Construction of asymptotically good, low-rate error-correcting codes through
  pseudo-random graphs.
\newblock {\em IEEE Trans. Inform. Theory}, 38(2):509--516, 1992.
\newblock [\epfmtdoi{10.1109/18.119713}]

\bibitem{AC}\bibhead{AC}
{\sc Noga Alon and Fan R.~K. Chung}: Explicit construction of linear sized
  tolerant networks.
\newblock {\em Discrete Math.}, 72(1--3):15--19, 1988.
\newblock Preliminary version in Proc. First Japan Conf. on Graph Theory and
  Appl., 1986.
\newblock [\epfmtdoi{10.1016/0012-365X(88)90189-6}]

\bibitem{AKRS}\bibhead{AKRS}
{\sc Noga Alon, Gil Kalai, Moty Ricklin, and Larry Stockmeyer}: Lower bounds on
  the competitive ratio for mobile user tracking and distributed job
  scheduling.
\newblock {\em Theoret. Comput. Sci.}, 130(1):175--201, 1994.
\newblock Preliminary version in
  \href{https://doi.org/10.1109/SFCS.1992.267757}{FOCS'92}.
\newblock [\epfmtdoi{10.1016/0304-3975(94)90158-9}]

\bibitem{AST}\bibhead{AST}
{\sc Noga Alon, Paul Seymour, and Robin Thomas}: A separator theorem for
  nonplanar graphs.
\newblock {\em J. Amer. Math. Soc.}, 3(4):801--808, 1990.
\newblock [\epfmtdoi{10.2307/1990903}]

\bibitem{Alweiss:19}\bibhead{Alweiss:19}
{\sc Ryan Alweiss}: Noisy corruption detection.
\newblock {\em Inform. Process. Lett.}, 155(105897), 2020.
\newblock [\epfmtdoi{10.1016/j.ipl.2019.105897}, \epfmt{arxiv}{1908.07493}]

\bibitem{BK99}\bibhead{BK99}
{\sc Piotr Berman and Marek Karpinski}: On some tighter inapproximability
  results (extended abstract).
\newblock In {\em Proc. 26th Internat. Colloq. on Automata, Languages and
  Programming (ICALP'99)}, volume 1644, pp. 200--209. Springer, 1999.
\newblock [\epfmtdoi{10.1007/3-540-48523-6\_17}]

\bibitem{CLRS}\bibhead{CLRS}
{\sc Thomas~H. Cormen, Charles~E. Leiserson, Ronald~L. Rivest, and Clifford
  Stein}: {\em Introduction to Algorithms}.
\newblock MIT Press, Cambridge, MA, third edition, 2009.

\bibitem{DPPU}\bibhead{DPPU}
{\sc Cynthia Dwork, David Peleg, Nicholas Pippenger, and Eli Upfal}: Fault
  tolerance in networks of bounded degree.
\newblock {\em SIAM J. Comput.}, 17(5):975--988, 1988.
\newblock Preliminary version in
  \href{https://doi.org/10.1145/12130.12169}{STOC'86}.
\newblock [\epfmtdoi{10.1137/0217061}]

\bibitem{fjeldstad2003fighting}\bibhead{fjeldstad2003fighting}
{\sc Odd-Helge Fjeldstad}: Fighting fiscal corruption: lessons from the
  {Tanzania Revenue Authority}.
\newblock {\em Public Administration and Development}, 23(2):165--175, 2003.
\newblock [\epfmtdoi{10.1002/pad.278}]

\bibitem{GO}\bibhead{GO}
{\sc Juan~A. Garay and Rafail Ostrovsky}: Almost-everywhere secure computation.
\newblock In {\em Advances in Cryptology---{EUROCRYPT'08}}, volume 4965, pp.
  307--323. Springer, 2008.
\newblock [\epfmtdoi{10.1007/978-3-540-78967-3\_18}]

\bibitem{HA74}\bibhead{HA74}
{\sc S.~Louis Hakimi and Ashok~T. Amin}: Characterization of connection
  assignment of diagnosable systems.
\newblock {\em IEEE Trans. Computers}, C-23(1):86--88, 1974.
\newblock [\epfmtdoi{10.1109/t-c.1974.223782}]

\bibitem{HLW}\bibhead{HLW}
{\sc Shlomo Hoory, Nathan Linial, and Avi Wigderson}: Expander graphs and their
  applications.
\newblock {\em Bull. Amer. Math. Soc.}, 43(4):439--561, 2006.
\newblock [\epfmtdoi{10.1090/S0273-0979-06-01126-8}]

\bibitem{JiMoRa:18}\bibhead{JiMoRa:18}
{\sc Yan Jin, Elchanan Mossel, and Govind Ramnarayan}: Being corrupt requires
  being clever, but detecting corruption doesn't.
\newblock In {\em 10th Innovations in Theoretical Comp. Sci. Conf. (ITCS'19},
  volume 124, pp. 45:1--45:14. Schloss Dagstuhl--Leibniz-Zentrum fuer
  Informatik, 2019.
\newblock [\epfmtdoi{10.4230/LIPIcs.ITCS.2019.45}, \epfmt{arxiv}{1809.10325}]

\bibitem{KTA75}\bibhead{KTA75}
{\sc Tiko Kameda, Shunichi Toida, and F.~J. Allan}: A diagnosing algorithm for
  networks.
\newblock {\em Information and Control}, 29(2):141--148, 1975.
\newblock [\epfmtdoi{10.1016/S0019-9958(75)90508-2}]

\bibitem{KR80}\bibhead{KR80}
{\sc Jon~G. Kuhl and Sudhakar~M. Reddy}: Distributed fault-tolerance for large
  multiprocessor systems.
\newblock In {\em Proc. 7th Ann. Symp. on Computer Architecture (ISCA'80)}, pp.
  23--30. ACM Press, 1980.
\newblock [\epfmtdoi{10.1145/800053.801905}]

\bibitem{LSP}\bibhead{LSP}
{\sc Leslie Lamport, Robert Shostak, and Marshall Pease}: The {B}yzantine
  generals problem.
\newblock {\em ACM Trans. Program. Lang. Syst.}, 4(3):382--401, 1982.
\newblock [\epfmtdoi{10.1145/357172.357176}]

\bibitem{LT}\bibhead{LT}
{\sc Richard~J. Lipton and Robert~Endre Tarjan}: A separator theorem for planar
  graphs.
\newblock {\em SIAM J. Appl. Math.}, 36(2):177--189, 1979.
\newblock [\epfmtdoi{10.1137/0136016}]

\bibitem{LPS}\bibhead{LPS}
{\sc Alex Lubotzky, Ralph Phillips, and Peter Sarnak}: Ramanujan graphs.
\newblock {\em Combinatorica}, 8(3):261--277, 1988.
\newblock [\epfmtdoi{10.1007/BF02126799}]

\bibitem{Ma}\bibhead{Ma}
{\sc Grigory~A. Margulis}: Explicit group-theoretic constructions of
  combinatorial schemes and their applications in the construction of expanders
  and concentrators ({R}ussian).
\newblock {\em Problemy Peredachi Informatsii}, 24(1):51--60, 1988.
\newblock
  \href{http://www.mathnet.ru/php/archive.phtml?wshow=paper&jrnid=ppi&paperid=686&option_lang=eng}{Link
  at Mathnet.ru.} English translation: Problems Inform. Transmission 24
  (1988/1) 39--46.

\bibitem{mukherjee1994network}\bibhead{mukherjee1994network}
{\sc Biswanath Mukherjee, L.~Todd Heberlein, and Karl~N. Levitt}: Network
  intrusion detection.
\newblock {\em IEEE Network}, 8(3):26--41, 1994.
\newblock [\epfmtdoi{10.1109/65.283931}]

\bibitem{nielsen2003corruption}\bibhead{nielsen2003corruption}
{\sc Richard~P. Nielsen}: Corruption networks and implications for ethical
  corruption reform.
\newblock {\em Journal of Business Ethics}, 42(2):125--149, 2003.
\newblock [\epfmtdoi{10.1023/A:1021969204875}]

\bibitem{PMC67}\bibhead{PMC67}
{\sc Franco~P. Preparata, Gernot Metze, and Robert~T. Chien}: On the connection
  assignment problem of diagnosable systems.
\newblock {\em IEEE Trans. Electronic Computers}, EC-16(6):848--854, 1967.
\newblock [\epfmtdoi{10.1109/PGEC.1967.264748}]

\bibitem{RS}\bibhead{RS}
{\sc Prasad Raghavendra and David Steurer}: Graph expansion and the unique
  games conjecture.
\newblock In {\em Proc. 42nd STOC}, pp. 755--764. ACM Press, 2010.
\newblock [\epfmtdoi{10.1145/1806689.1806792}]

\bibitem{rock2004comparative}\bibhead{rock2004comparative}
{\sc Michael~T. Rock and Heidi Bonnett}: The comparative politics of
  corruption: accounting for the {E}ast {A}sian paradox in empirical studies of
  corruption, growth and investment.
\newblock {\em World Development}, 32(6):999--1017, 2004.
\newblock [\epfmtdoi{10.1016/j.worlddev.2003.12.002}]

\bibitem{lgorzata2004survey}\bibhead{lgorzata2004survey}
{\sc Ma{\l}gorzata Steinder and Adarshpal~S. Sethi}: A survey of fault
  localization techniques in computer networks.
\newblock {\em Sci. Comput. Programming}, 53(2):165--194, 2004.
\newblock [\epfmtdoi{10.1016/j.scico.2004.01.010}]

\bibitem{tanner}\bibhead{tanner}
{\sc R.~Michael Tanner}: Explicit construction of concentrators from
  generalized $n$-gons.
\newblock {\em SIAM J. on Algebraic Discrete Methods}, 5(3):287--293, 1984.
\newblock [\epfmtdoi{10.1137/0605030}]

\end{thebibliography}

\end{document}
